\section{总结与延伸}

\subsection{核心结论}

\begin{enumerate}
\item 小鹏转向 Physical AI，不是给车加大模型，而是重写自动驾驶的数据和模型闭环。
\item 刘先明的路径说明，真实世界问题和使命感会持续牵引技术选择。
\item 自动驾驶软件栈从规则、半模型、端到端，走向更大模型和云端工厂。
\item “Language 是毒药”的准确含义，是语言作为中间监督瓶颈会限制 data scaling。
\item 拆掉中间层的目标，是让连续传感器数据和动作输出更直接地进入训练。
\item 主机厂优势在真实车队、可控数据链路和量产反馈。
\item 云端模型工厂负责训练大模型，再蒸馏、量化、剪枝到车端部署。
\item 技术简化和组织扁平化是同一件事：都在减少瓶颈。
\item 换帅背后是技术阶段切换，刘先明的任务是把 Physical AI 路线跑成生活服务。
\item 未来关键 bet 是 Physical AI 中的 scaling，以及持续识别自己的思维误区。
\end{enumerate}

\subsection{开放问题}

最后保留开放问题，是因为小鹏这条路线仍在高速展开。真正要观察的，不只是“拆掉了什么”，还包括拆掉之后系统是否更稳、更可控、更能 scale。

\begin{itemize}
\item 拆掉 Language 后，模型如何保留足够高层语义和可解释性？
\item 云端大模型到车端模型的蒸馏，会不会成为体验上限？
\item 主机厂数据闭环能否追上 Robotaxi 专用车队的高质量闭环？
\item Robotaxi 在广州跑好，距离全国规模化还有哪些非技术阻力？
\item 扁平工程组织在规模化量产和质量体系中能保持多久？
\item Physical AI 的 scaling 会先在车上验证，还是在机器人/其他本体上验证？
\end{itemize}

\subsection{拓展阅读}

\begin{itemize}
\item EP121 谭杰访谈：机器人、世界模型、跨本体和 Gemini Robotics 1.5。
\item EP132 高继扬访谈：Waymo、Momenta 和星海图的具身智能生产化。
\item EP134 数据综述：Data Recipe、机器人数据和数据定价。
\item Software 2.0、VLA、VLM、自监督学习和模型蒸馏相关资料。
\end{itemize}

\begin{importantbox}{最后的判断}
EP120 最值得保留的不是一句“拆掉 L”，而是一套工程哲学：当中间层成为数据 scaling 的瓶颈，就要拆；当组织层级成为问题反馈的瓶颈，也要拆。Physical AI 的进步，来自更直接的数据、更大的模型、更短的工程链路和更坚决的量产验证。
\end{importantbox}

\end{document}
